% ────────────────────  English abstract  ────────────────────
\cleardoublepage
\pagenumbering{roman}
\setcounter{page}{2}
\phantomsection              % for hyperlinks/bookmarks
\chapter*{{\Huge\bfseries Abstract}}   % big on the page
\thispagestyle{noheader}
\addcontentsline{toc}{chapter}{Abstract} % normal in TOC

\vspace{0.5cm}


Wheat is a globally essential cereal crop, but its yield is significantly threatened by a variety of diseases. Traditional diagnostic methods, which rely on manual visual inspection, are often time-consuming, subjective, and inaccessible to many farmers, especially in remote areas. With the emergence of smart agriculture and the advancement of deep learning, automated image-based disease recognition has become a promising solution.

This project aims to design and implement a robust system based on convolutional neural networks (CNNs) for the detection and classification of wheat diseases using images of various plant parts such as leaves, stems, and spikes. Our methodology involves a comprehensive literature review, data collection and preprocessing, model training and evaluation, and finally the integration of the solution into an interactive web platform tailored for farmers. To enhance system utility, we also include insect pest detection and an AI assistant module accessible via API to guide users.

We propose a hierarchical classification framework composed of two stages. In the first stage, multiple CNN architectures are individually combined with a Support Vector Machine (SVM) classifier to distinguish between healthy and diseased samples, forming a hybrid CNN-SVM approach. The second stage employs a ResNet50 backbone enhanced with an Efficient Channel Attention (ECA) module to identify nine specific wheat diseases. Experimental results show that the binary classifier (Model 1) achieved an accuracy of 99.60\%, while the multi-class classifier (Model 2) reached 99.83\% accuracy. The full system attained an overall accuracy of 99.48\%.

The proposed solution contributes to smart agriculture by enabling early, accessible, and reliable diagnosis, thus aiding in timely intervention and improving overall crop management.


\vspace{1cm}
    \noindent\rule{\linewidth}{1pt}
    \noindent
    \textbf{Keywords —} Smart Agriculture, Deep Learning, Convolutional Neural Networks, Wheat Disease Detection, Computer Vision.\\
    \noindent\rule{\linewidth}{1pt}






 

% ────────────────────  French résumé  ────────────────────
\cleardoublepage
\pagenumbering{roman}
\setcounter{page}{2}
\chapter*{{\Huge\bfseries Résumé}}
\thispagestyle{noheader}
\addcontentsline{toc}{chapter}{Résumé}

\vspace{0.5cm}

Le blé est une céréale essentielle à l’échelle mondiale, mais son rendement est fortement menacé par diverses maladies. Les méthodes de diagnostic traditionnelles, basées sur l’inspection visuelle manuelle, sont souvent chronophages, subjectives et inaccessibles pour de nombreux agriculteurs, notamment dans les zones éloignées. Avec l’émergence de l’agriculture intelligente et les avancées de l’apprentissage profond, la reconnaissance automatisée des maladies à partir d’images constitue une solution prometteuse.

Ce projet vise à concevoir et mettre en œuvre un système robuste basé sur des réseaux de neurones convolutifs (CNN) pour la détection et la classification des maladies du blé à partir d’images de différentes parties de la plante telles que les feuilles, les tiges et les épis. Notre méthodologie comprend une revue de littérature approfondie, la collecte et le prétraitement des données, l’entraînement et l’évaluation des modèles, et enfin l’intégration de la solution dans une plateforme web interactive destinée aux agriculteurs. Afin d’augmenter l’utilité du système, nous avons également intégré un module de détection des insectes nuisibles et un assistant intelligent accessible via une API pour guider les utilisateurs.

Nous proposons une approche hiérarchique de classification en deux étapes. Dans la première, plusieurs architectures CNN sont combinées individuellement avec un classificateur à vecteurs de support (SVM) pour distinguer les échantillons sains des échantillons malades, formant ainsi une approche hybride CNN-SVM. La seconde étape utilise l’architecture ResNet50 enrichie par un module d’attention efficace aux canaux (ECA) pour identifier neuf maladies spécifiques du blé. Les résultats expérimentaux montrent que le classificateur binaire (Modèle 1) a atteint une précision de 99{,}60\%, tandis que le classificateur multi-classes (Modèle 2) a atteint une précision de 99{,}83\%. Le système complet a obtenu une précision globale de 99{,}48\%.

La solution proposée contribue à l’agriculture intelligente en permettant un diagnostic précoce, accessible et fiable, facilitant ainsi une intervention rapide et une meilleure gestion des cultures.

\vspace{1cm}

    \noindent\rule{\linewidth}{1pt}
    \noindent
    \textbf{Mots clés —} Agriculture intelligente, Apprentissage profond, Réseaux de neurones convolutifs, Détection des maladies du blé, Vision par ordinateur.\\
    \noindent\rule{\linewidth}{1pt}






% ────────────────────  Arabic ملخص  ────────────────────

\cleardoublepage
\pagenumbering{roman}
\setcounter{page}{2}
\phantomsection              % for hyperlinks/bookmarks
\chapter*{\Huge\bfseries\hfill\textRL{مـلـخـص}}
\thispagestyle{noheader}
\addcontentsline{toc}{chapter}{\textRL{مـلـخـص}}


\vspace{0.1cm}

\begin{RLtext}

يُعد القمح من المحاصيل الزراعية الأساسية على مستوى العالم، إلا أن إنتاجه يواجه تهديدًا كبيرًا بسبب تنوع الأمراض التي تصيبه. تُعتبر طرق التشخيص التقليدية، المعتمدة على الفحص البصري اليدوي، بطيئة، وذات طابع ذاتي، وغالبًا ما تكون غير متاحة للمزارعين، خاصة في المناطق النائية. ومع ظهور الزراعة الذكية وتطور تقنيات التعلم العميق، أصبحت عملية التعرف التلقائي على الأمراض من خلال الصور حلاً واعدًا.

يهدف هذا المشروع إلى تصميم وتنفيذ نظام فعّال يعتمد على الشبكات العصبية الالتفافية (CNN) للكشف عن أمراض القمح وتصنيفها باستخدام صور لمختلف أجزاء النبات مثل الأوراق والسيقان والسُنابل. تشمل منهجيتنا مراجعة شاملة للدراسات السابقة، وجمع البيانات ومعالجتها، وتدريب النماذج وتقييمها، ثم دمج الحل في منصة ويب تفاعلية موجهة للمزارعين. ولتعزيز فاعلية النظام، قمنا أيضًا بإدراج وحدة لاكتشاف الحشرات الضارة ومساعد ذكي يعمل عبر واجهة برمجة التطبيقات (API) لتوجيه المستخدمين.

نقترح إطارًا هرميًا للتصنيف يتكون من مرحلتين: في المرحلة الأولى، يتم دمج عدة هياكل CNN بشكل فردي مع خوارزمية آلة متجه الدعم (SVM) للتمييز بين العينات السليمة والمصابة، مشكلين بذلك مقاربة هجينة CNN-SVM. أما في المرحلة الثانية، فنستخدم بنية ResNet50 معززة بوحدة الانتباه الكفء للقنوات (ECA) للتعرف على تسع أمراض محددة تصيب القمح. أظهرت النتائج التجريبية أن المصنف الثنائي حقق دقة بلغت 99{،}60\%، بينما بلغ المصنف متعدد الأصناف دقة قدرها 99{،}83\%. وبلغت الدقة الإجمالية للنظام الكامل 99{،}48\%.

تُساهم الحلول المقترحة في تطوير الزراعة الذكية من خلال تمكين التشخيص المبكر والموثوق، مما يساعد في التدخل في الوقت المناسب وتحسين إدارة المحاصيل الزراعية.

\end{RLtext}





\vspace{0.1cm}
\hspace*{0mm}\rule{\linewidth}{1pt}
\vspace{0.1cm}
\begin{RLtext}
\noindent\textbf{الكلمات المفتاحية —} الزراعة الذكية، التعلُّم العميق، الشبكات العصبية الالتفافية، الرؤية الحاسوبية، اكتشاف أمراض القمح.
\end{RLtext}
\vspace{0.1cm}
\rule{\linewidth}{1pt}






% \cleardoublepage
% \pagenumbering{roman}
% \setcounter{page}{2}
% \phantomsection  
% \begin{flushright}
%     {\Huge\bfseries\RL{مـلـخـص}}
% \end{flushright}
% \addcontentsline{toc}{chapter}{\textRL{مـلـخـص}}




% \begin{RLtext}

% يُعد القمح من المحاصيل الزراعية الأساسية على مستوى العالم، إلا أن إنتاجه يواجه تهديدًا كبيرًا بسبب تنوع الأمراض التي تصيبه. تُعتبر طرق التشخيص التقليدية، المعتمدة على الفحص البصري اليدوي، بطيئة، وذات طابع ذاتي، وغالبًا ما تكون غير متاحة للمزارعين، خاصة في المناطق النائية. ومع ظهور الزراعة الذكية وتطور تقنيات التعلم العميق، أصبحت عملية التعرف التلقائي على الأمراض من خلال الصور حلاً واعدًا.

% يهدف هذا المشروع إلى تصميم وتنفيذ نظام فعّال يعتمد على الشبكات العصبية الالتفافية (CNN) للكشف عن أمراض القمح وتصنيفها باستخدام صور لمختلف أجزاء النبات مثل الأوراق والسيقان والسُنابل. تشمل منهجيتنا مراجعة شاملة للدراسات السابقة، وجمع البيانات ومعالجتها، وتدريب النماذج وتقييمها، ثم دمج الحل في منصة ويب تفاعلية موجهة للمزارعين. ولتعزيز فاعلية النظام، قمنا أيضًا بإدراج وحدة لاكتشاف الحشرات الضارة ومساعد ذكي يعمل عبر واجهة برمجة التطبيقات (API) لتوجيه المستخدمين.

% نقترح إطارًا هرميًا للتصنيف يتكون من مرحلتين: في المرحلة الأولى، يتم دمج عدة هياكل CNN بشكل فردي مع خوارزمية آلة متجه الدعم (SVM) للتمييز بين العينات السليمة والمصابة، مشكلين بذلك مقاربة هجينة CNN-SVM. أما في المرحلة الثانية، فنستخدم بنية ResNet50 معززة بوحدة الانتباه الكفء للقنوات (ECA) للتعرف على تسع أمراض محددة تصيب القمح. أظهرت النتائج التجريبية أن المصنف الثنائي (النموذج 1) حقق دقة بلغت 99{،}60\%، بينما بلغ المصنف متعدد الأصناف (النموذج 2) دقة قدرها 99{،}83\%. وبلغت الدقة الإجمالية للنظام الكامل 99{،}48\%.

% تُساهم الحلول المقترحة في تطوير الزراعة الذكية من خلال تمكين التشخيص المبكر والموثوق، مما يساعد في التدخل في الوقت المناسب وتحسين إدارة المحاصيل الزراعية.

% \end{RLtext}